\chapter[Stan and Walnutpie]{Stan and Walnutpie: Source Transfer, XLA Execution, and Evidence}
\label{ch:stan-walnutpie}
\begingroup
\emergencystretch=2em
% Permit long source paths to wrap locally without changing other chapters.
\expandafter\def\expandafter\UrlBreaks\expandafter{\UrlBreaks
\do\a\do\b\do\c\do\d\do\e\do\f\do\g\do\h\do\i\do\j
\do\k\do\l\do\m\do\n\do\o\do\p\do\q\do\r\do\s\do\t
\do\u\do\v\do\w\do\x\do\y\do\z\do\-\do\_}

This chapter records the standalone sampler programme completed between
27 and 31 October 2026. Its motivating question was practical: can an
adaptive HMC implementation accept a TensorFlow filtering model, compile
the repeated numerical work with XLA, and approach the execution efficiency
of native Stan? The subsequent experiments also asked whether local
step-size refinement helps difficult geometry, whether CPU tree control
should be combined with GPU model evaluation, and whether Walnutpie can
recover a saved financial posterior through a frozen NeuTra transform.

These questions have different answers. The supported Stan algorithm and
Walnutpie mechanisms were transferred and checked against pinned sources.
Native C++ calls can be embedded in an XLA graph, but that does not move
their computation onto the GPU. The separate TensorFlow implementations do
evaluate the model and sampler numerics on the selected XLA device. Cheap
analytic targets often favour CPU execution; expensive batched models can
favour GPU execution. The full80 financial comparison passed its declared
posterior checks, whereas the separate q20 application did not. Neither
source agreement nor successful compilation alone establishes posterior
accuracy or a performance advantage.

The scope is the owner-authorized standalone work in
\path{docs/benchmarks/}. It does not restore excluded public NUTS routes,
replace the repository's public HMC tuning procedures, or make these
experimental interfaces new production defaults. Historical and failed
attempts remain part of the evidence.

\section{Which implementation is being compared?}
\label{sec:sw-implementations}

Stan 2.37.0 is the native reference throughout this programme, not a claim
about the latest Stan release. Its source was installed locally from a
checksummed release archive. The relevant authority is the C++ Euclidean
multinomial NUTS implementation and its service-level initialization and
adaptation, rather than the original slice-based NUTS pseudocode alone
\citep{swStan237,swNUTS2014}. TensorFlow 2.20.0 and TFP 0.25.0 were the
tested framework versions. The following implementations must be kept distinct.

\begin{description}[style=nextline]
\item[Stan-style warmup around TFP NUTS.]
\path{stan_nuts_warmup.py} transfers step-size search, dual averaging and
windowed diagonal/dense metric learning, while TFP supplies the transition.
Its arithmetic helpers are reusable TensorFlow functions; its controller
is specifically wired to TFP NUTS. It is not Stan's tree implementation.
\item[Direct Stan algorithm port.]
\path{stan_nuts.py} translates the pinned Euclidean multinomial tree,
integrator, metric, adaptation and flat-coordinate initialization into
TensorFlow. Native Stan is an oracle and comparator, not a runtime dependency.
\item[Native TFP interface.]
The first \path{stan_nuts_tfp.py} wrapped the direct tree in TFP's
\texttt{TransitionKernel} protocol. On 28 October it was replaced by stock
TFP NUTS and \texttt{sample\_chain}, with the checked Stan-style warmup and
fixed-metric transformations. An API wrapper alone did not accelerate the
direct tree; the rewrite changed the numerical implementation.
\item[Native C++ XLA bridge.]
\path{native_sampler_bridge/} calls native Stan or Walnutpie C++ from an
XLA custom call. Its GPU entry stages inputs to host, executes the CPU
sampler, and returns outputs to device. This preserves native computation
but is not GPU sampler execution.
\item[TensorFlow models inside GPU/XLA samplers.]
\path{stan_gpu_tf/} reuses the direct Stan authority with explicit model
tensor operands; \path{walnutpie_gpu_tf/} implements Walnutpie numerics
in TensorFlow. These paths contain the model, its position score and the
sampler loops in compiled functions, without a CPU sampler callback.
\end{description}

All paths above are relative to \path{docs/benchmarks/}. Transferring this
supported sampler is not transferring the whole Stan compiler, constrained
parameter language, Math library, CSV/service conventions, RNG stream or
all HMC variants. In particular, a C++ callback cannot generally invoke an
arbitrary TensorFlow graph and thereby become fused GPU computation.

\section{Stan adaptation: mathematics, order, and audit}
\label{sec:sw-stan-tuning}

Let $q\in\mathbb R^d$, $U(q)=-\log\pi(q)$, and let $C$ denote the
inverse momentum covariance. The convention used by the direct Stan port is
\begin{equation}
 H(q,p)=U(q)+\tfrac12p^\top Cp,\qquad
 p\sim N(0,C^{-1}),\qquad \dot q=Cp.
 \label{eq:sw-hamiltonian}
\end{equation}
Unit, diagonal and dense positive-definite $C$ are supported. Momentum
generation, velocity and kinetic energy must use the same convention.
If $C=AA^\top$, $q=Az$ and $p=A^{-\top}p_z$, then
$p^\top dq=p_z^\top dz$ and the kinetic energy is $p_z^\top p_z/2$.
This gives the constant-metric coordinate change used around unit-mass
TFP NUTS. The constant log determinant does not affect the fixed-metric
transition; at a metric update the physical position must be preserved,
the new $z$ solved for, and density/gradient caches rebuilt. Reusing a cache
from the old coordinates would be wrong.

\subsection{Three warmup stages}

Stan's initial fast interval learns step size, expanding slow intervals
learn metric and step size, and a terminal fast interval learns step size
under the final metric. The inherited buffers are 75 initial, 25 in the
first slow window and 50 terminal transitions. For 1000 warmup transitions,
the zero-based metric-update indices are $99,149,249,449,949$. These are
source-derived starting choices, not universal adequacy claims.

Within a slow window, Welford updates estimate its covariance $S_n$. For
$n>1$ the regularized inverse metric is
\begin{equation}
 C_n=\frac{n}{n+5}S_n+\frac{5}{n+5}10^{-3}I.
 \label{eq:sw-metric-shrinkage}
\end{equation}
The diagonal route uses only the diagonal. A one-observation window retains
the previous covariance estimate before the same shrinkage, following the
pinned source. The accumulator resets at a window boundary; it does not
pool every warmup observation. A fresh reasonable-step search follows the
metric installation, then dual averaging restarts in the new geometry.

For target acceptance $\delta$, observed trajectory-average Metropolis
probability $\alpha_t$, and shrinkage centre $\mu$, the transferred
dual-averaging recursion is
\begin{align}
 \bar H_t&=\left(1-\frac1{t+t_0}\right)\bar H_{t-1}
             +\frac{\delta-\alpha_t}{t+t_0},\label{eq:sw-dual-h}\\
 \log\epsilon_t&=\mu-\frac{\sqrt t}{\gamma}\bar H_t,\\
 \log\bar\epsilon_t&=t^{-\kappa}\log\epsilon_t
                   +(1-t^{-\kappa})\log\bar\epsilon_{t-1}.
\end{align}
The source defaults are $\delta=.8$, $\gamma=.05$, $\kappa=.75$ and
$t_0=10$. Retained sampling freezes the final smoothed step size and
metric. The statistic $\alpha_t$ averages over trajectory leapfrogs;
it is not the fraction of transitions whose selected position changed.
Rejected nonfinite leaves contribute zero acceptance while a valid retained
state remains usable.

Initial and window-boundary ordering matters. The initial service sets
$\mu=\log(10\epsilon_{\rm supplied})$ before its first search; at a
metric boundary the centre uses the newly searched step. The search
brackets a one-leapfrog acceptance probability of $.8$ by doubling or
halving with the source's momentum-refresh and repeated-probe cadence.
Its $.8$ bracket is separate from configurable $\delta$. Short warmup
also has pinned-version semantics: fewer than 20 transitions do not learn
a metric, and the inspected 2.37 fallback for shortened buffers returns
without restarting the fresh window counter. The port preserves that
behaviour instead of silently replacing it with a more convenient schedule.

\subsection{What the tests establish}

The adaptation audit compiled the actual Stan headers, rather than testing
against another transcription. It checked covariance and singleton-window
arithmetic, dual-averaging updates/restarts, initial search ties and failures,
and native adapter replay at each metric boundary. Dense-coordinate
orientation was checked against the native leapfrog integrator. State,
density and score caches were independently recomputed around transitions;
continuation, per-chain separation, invalid inputs and final freezing were
covered. The reported audit suite passed 64 tests with zero skips; CPU and
GPU XLA execution was separately qualified.

The direct port additionally checks recursive tree behaviour through
supplied momenta and uniform random tapes: proposal selection, early return,
work/acceptance counters, strict energy thresholds, direction ties, the
three momentum-based U-turn checks and RNG consumption. Sources include
\path{stan/src/stan/mcmc/hmc/nuts/base_nuts.hpp},
\path{mcmc/hmc/base_hmc.hpp}, \path{mcmc/stepsize_adaptation.hpp},
the windowed variance/covariance adapters and service initialization.
The combined direct-port record reports 128 tests. These fixtures establish
the checked source behaviour, not universal correctness over all targets.

The generic target returns FP64 log densities of shape $[C]$ from positions
of shape $[C,d]$, with independent rows. It owns constraints and Jacobians.
Flat supplied/partial/random starts, initialization retries, jitter,
adaptation-off, zero warmup, thinning and retained continuation are supported
within the documented interface. TF Philox keys/counters replace Stan's
random engine: identical numeric seeds do not imply identical trajectories.
Arbitrary Python exceptions cannot run inside XLA, so compiled validity
flags must be checked at host boundaries. The reusable guides are
\path{docs/reference/nuts-tuning.md}, \path{nuts-direct-stan.md},
\path{nuts-stan-tfp.md} and \path{nuts-stan-coverage.md} in that directory.

\section{Walnutpie and the WALNUTS paper}
\label{sec:sw-walnuts}

The deterministic C++ authority is Walnutpie commit
\texttt{5710b0334ba6}; the paper experiments
use the separate author WALNUTS repository at
\texttt{895a9b7a595b}
\citep{swWalnutpieSource,swWalnutsAuthor}. Full commit hashes appear in the bibliography and source manifests.
These are related sources, not interchangeable implementations. The published account is
\citet{swWalnuts2026}, particularly Sections 3.1--3.4, 4.1, 4.3,
Remark 13, equations (13)--(14), and Appendices A--C.

\subsection{Local refinement and reverse recoverability}

A macro interval of length $h$ is subdivided into microsteps. In the local
deterministic code, with minimum count $m$ and refinement exponent $k$,
there are $m2^k$ leapfrogs of size $h/2^k$; the macro duration is therefore
$mh$. Each candidate is accepted locally only after its energy test and
coarser reverse-trajectory checks. Merely shrinking the step in a difficult
region without checking reverse selection does not reproduce this sampler.
The optimized source uses absolute endpoint energy error; the paper's
Remark 13 explicitly permits that deterministic selector in place of the
more conservative micro-trajectory energy envelope. Intermediate nonfinite
states must still be handled.

For randomized refinement, write $r_f$ and $r_b$ for the probabilities of
selecting the realised microstep count from the forward and reverse endpoints.
The paper's forward orbit-weight recursion can be written
\begin{equation}
 \log w_j=\log w_{j-1}-H_j+H_{j-1}+\log r_b-\log r_f.
 \label{eq:sw-orbit-weight}
\end{equation}
The negative-index recursion has the corresponding directional probabilities.
These cumulative factors are necessary: state-dependent integration is not
just ordinary fixed-step HMC with a locally convenient epsilon. Zero reverse
support gives zero weight. The implementation keeps the orbit geometry while
representing zero probability by $-\infty$ in log space.

In the R2P implementation, the critical refinement exponent $k_*$ is selected
by a bounded energy search; exponent $k_*$ is chosen with probability $2/3$,
and $k_*+1$ with probability $1/3$. The paper experiments also use $.2$
jitter. The capped author selector can return its cap without satisfying
the threshold, and the random finer choice may then use cap plus one.
The port exposes this distinction rather than claiming the bound certifies
all energy tests.

An audit found a missing second cumulative backward log-weight update in
the pinned author Python tree near line 462. The local R2P tree includes
the term required by equations (13)--(14). Its oracle therefore combines
the unmodified author macro with an independently written recursive tree
using explicit orbit points and the paper equations. A negative control
removes the disputed term. Literal original-author and corrected-author
chains remain separate experiment arms. This is a documented source
difference, not a claim of bitwise agreement with the original whole chain.

\subsection{Adaptation and run control}

Deterministic Walnutpie uses displacement U-turns, Barker selection within
subtrees and biased Metropolis selection when extending the outer tree.
Its base-macro quality statistic is $\exp(-|\Delta H|)$; the default
target $.8$ feeds Adam at base-macro attempts. Discounted draw/score moments
adapt the diagonal mass, and average macro work adapts the minimum microstep
count. It does not substitute Stan's dual averaging or covariance windows.
A fixed-HMC endpoint acceptance target such as $.70$ cannot be copied over
as if it controlled the same quantity.

The TensorFlow implementation adds supplied, zero and normal starts;
supplied, unit and absolute-score diagonal masses; optional geometric mass
averaging; supplied steps or source-style fixed-momentum initialization
search. The default 64-update search safety bound can explicitly be removed;
an uncapped compiled search has no internal wall-clock interruption.
Invalid positions, steps, reciprocal masses and duplicate stream keys
are rejected through validity outputs.

Compiled controllers reproduce the pinned relative step/mass agreement
statistics and legacy density-agreement statistic. Synchronized batched
advancement replaces asynchronous C++ workers. Host deadlines, callbacks
and cancellation are observed between compiled chunks. Thus a timeout
boundary may include an in-flight compile or transition. Full checkpoints
preserve random state, caches, Adam state, online moments and iteration
counts; continuing retained sampling must not restart adaptation. A
maximum-iteration stop or agreement of tuning parameters is not posterior
convergence. Appendix C's paper quantile-tuning experiment remains a
separate CPU reference procedure; optional Adam warmup around R2P does
not claim to reproduce it.

\subsection{Supplied mass and NeuTra}

Walnutpie's supplied API takes momentum covariance $M$, whereas the Stan
port takes its inverse $C$. Fixed mode retains $M$ exactly; adaptive
diagonal mode seeds the existing draw/score estimator from it. For a dense
supplied $M=LL^\top$, the implementation uses
\begin{equation}
 y=L^\top q,\qquad p=Lr,\qquad
 \log\pi_y(y)=\log\pi(L^{-\top}y)-\log|\det L|.
 \label{eq:sw-dense-mass}
\end{equation}
Substitution gives $p^\top dq=r^\top dy$ and
$p^\top M^{-1}p=r^\top r$. Adaptive dense-interface mode learns
$L\operatorname{diag}(m)L^\top$ in this fixed basis, not an unrestricted
dense covariance. The physical displacement U-turn rule must also be
transformed; naive reuse of a Euclidean $y$-space rule would change the
specified algorithm. With adaptive inverse diagonal entries $D$ in the
$y$ basis, the physical dot product becomes
$r^\top D L^{-1}L^{-\top}\Delta y$.

For a frozen learned transform $q=T(z)$, the latent target is
\begin{equation}
 \log\pi_z(z)=\log\pi_q(T(z))+\log|\det J_T(z)|.
 \label{eq:sw-neutra-target}
\end{equation}
The requested full80 experiment fixes $M_z=I$ and adapts only the permitted
step/work controls. Learning another ordinary metric is not a prerequisite.
The Jacobian must appear exactly once, and the cached score must differentiate
this same transformed density.

\section{What full XLA execution means here}
\label{sec:sw-xla}

Both TensorFlow sampler families expose initialization, warmup, sampling
and value/gradient functions with stable \texttt{input\_signature}s
and \texttt{jit\_compile=True}. Model data are runtime tensor operands;
changes in values with the same signature need not retrace. Shape or static
configuration changes require a new compiled function. Independent-row
batch semantics are checked: cross-chain normalization, mutable targets,
random likelihood values or missing Jacobians violate the interface.
Numerical loops use native TensorFlow control flow, without pfor or Python
sample-by-sample target evaluation.

The native bridge instead lowers a \texttt{BfNativeSamplerV1} custom call.
GPU producer/bridge/consumer tests verify that staging composes correctly
inside a GPU graph, but the sampler remains on CPU. Its fixed analytic
target library is not a general fusion interface for arbitrary TensorFlow
models. Native direct and CPU/GPU-bridge outputs were byte-identical on
the tested fits; same-engine overhead is measured separately from
cross-algorithm efficiency.

For the genuine GPU implementations, optimized HLO, device output placement,
runtime-data sensitivity, one-trace checks, absence of forbidden callbacks
and device profiling establish the tested numerical execution. This does
not imply one fused CUDA kernel, zero host scheduling, or gradients through
the accept/reject sampler itself. Device BLAS/solver custom calls are not
CPU sampler callbacks. TensorFlow's command-buffer and profiling modes may
change scheduling: the inspected TF 2.20 profiler disables command buffers.
Reported first-call times include tracing, compilation and an execution,
unless a compiler-only alarm is explicitly quoted. Warmup and retained
functions can incur separate cold compilations.

The supplied GPU launchers configure and verify on-demand memory growth
before initialization. CPU reference processes explicitly hide GPUs.
All timing tables below identify their workload and boundary; GPU allocator
peak, host RSS and driver reservation are different quantities.

\section[Stan comparisons]{From a banana benchmark to a source-matched comparison}
\label{sec:sw-stan-results}

\subsection{Cheap targets and hidden batched work}

The first pilot used FP64 TFP NUTS, identity mass, four chains, 128 warmup
and 256 retained transitions per chain, three seed/start repetitions and
maximum depth six. The larger product-banana case used 32 chains and
32 dimensions. Retained batch medians were as follows; warmup, cold
compilation and file output are excluded.

\begin{center}\small
\begin{tabular}{lrrr}
\toprule
Target & CPU seconds & GPU seconds & GPU/CPU\\
\midrule
Gaussian, $d=2$, four chains & .007285 & .090529 & 12.43\\
Banana, $d=2$, four chains & .013889 & .190009 & 13.68\\
Product banana, $d=32$, 32 chains & .202665 & .757104 & 3.74\\
\bottomrule
\end{tabular}
\end{center}

This nominated CPU/XLA for these small fixtures. Nonrandomized device order,
short chains and different adaptation paths prevent a universal CPU claim.
The subsequent finite-grid banana comparison selected fixed HMC $L=3$.
For 4096 retained transitions in four chains, its median was $.0273044$
seconds versus $.1405021$ for NUTS. The reported leapfrog counts were
$3$ versus $4.8323$, but the selected NUTS candidate had one stationary
chain, maximum R-hat $1.52671$ and an unavailable tail diagnostic. HMC's
maximum R-hat was $1.00136$. The NUTS candidate failed the validity screen;
the observed timing/ESS ratios cannot establish posterior efficiency.

An instrumented replay explained why ``only 1.83 extra leapfrogs'' was
misleading. TFP's batched loops continue while any chain remains active;
the integrator evaluates the dense batch before masking completed rows.
The following counters subtract bootstrap evaluation and were checked
against uninstrumented outputs; instrumented calls were not timed.

\begin{center}\small
\begin{tabular}{lrr}
\toprule
Method / chain batch & Reported steps per chain & Dense gradient calls per transition\\
\midrule
HMC / 4 & 3.000 & 3.000\\
NUTS / 1 & 4.782 & 4.782\\
NUTS / 4 & 4.869 & 9.395\\
NUTS / 16 & 4.920 & 16.532\\
\bottomrule
\end{tabular}
\end{center}

The counters identify additional executed work, not its complete cost.
Tree bookkeeping, selection, random generation, synchronization and memory
traffic remain. A broader identity-metric pilot tested Gaussian32, rotated
Gaussian32 with covariance condition 1000, banana8, centered funnel10,
an exact noncentered funnel control and regression512$\times$32. NUTS
passed the short rotated-Gaussian screen when the selected fixed-HMC
candidate did not; both failed the centered funnel; both passed the
noncentered control. Thus the initial cheap-banana result was not a reason
to reject adaptive trajectory lengths on harder geometry.

\subsection{Tuning-inclusive Stan comparisons}

The first Stan-style warmup comparison used installed CmdStan 2.37.0,
checked native target values/scores and separately adapted metrics. Both
engines passed all three seeds on five of six targets; both failed the
centered funnel. TF/native paired elapsed ratios on the passing targets
were $5.41$, $4.05$, $6.43$, $6.98$ and $2.71$ for Gaussian32, rotated
Gaussian32, banana8, noncentered funnel10 and regression512$\times$32,
respectively. Their reported approximate intervals excluded one, but the
later adaptation audit repaired edge cases. Those timings retain their
original source versions and are not measurements of the repaired code.

The direct-port comparison then used seven targets, three seeds, FP64,
four parallel scalar CPU chains, 1000 warmup and 2000 retained transitions
per chain. The dense metric was used for the rotated Gaussian. Reusable
TF kernels were already compiled. TF costs include initialization, search,
adaptation and sampling; Stan process times include launch and CSV output.
TF serialization and common diagnostics are recorded separately.

\begin{center}\small
\begin{tabular}{lrrrl}
\toprule
Target & Direct TF & Audited TFP & Native Stan & Passing seeds\\
 & \multicolumn{3}{c}{Median seconds} & per engine\\
\midrule
Gaussian32 & 1.047 & .628 & .123 & 3/3\\
Rotated Gaussian32 & 2.631 & 1.350 & .344 & 3/3\\
Banana8 & 1.038 & .624 & .100 & 3/3\\
Centered funnel10 & 2.151 & 1.068 & .231 & 0/3\\
Noncentered funnel10 & .635 & .516 & .074 & 3/3\\
Regression512$\times$32 & 1.293 & 1.010 & .344 & 3/3\\
Regression4096$\times$64 & 7.160 & 9.972 & 3.732 & 3/3\\
\bottomrule
\end{tabular}
\end{center}

The direct port reproduced the tested transition rules but did not recover
native Stan speed. Different RNGs and dynamic work also mean these totals
are not fixed-work compiler comparisons. The associated public fixed-HMC
procedure ran 21 fits and assessed all 216 verified candidate-set members;
13 members passed in seven fits. Some target cohorts had no passing member,
and one search exhausted its per-fit budget before finishing. This bounded
exercise includes preparation/search costs but is not the best achievable
fixed-HMC baseline or an isolated experiment on trajectory-length choice.

After the native-TFP rewrite, the actual \texttt{TransitionKernel}/
\texttt{sample\_chain} path was measured, rather than inferring its speed
from an older manual loop. The same seven-target, three-seed design gave:

\begin{center}\small
\begin{tabular}{lrrl}
\toprule
Target & Native TFP seconds & Stan seconds & Passing seeds each\\
\midrule
Gaussian32 & .610 & .121 & 3/3\\
Rotated Gaussian32 & 1.355 & .330 & 3/3\\
Banana8 & .629 & .095 & 3/3\\
Centered funnel10 & 1.082 & .168 & 0/3\\
Noncentered funnel10 & .506 & .074 & 3/3\\
Regression512$\times$32 & .972 & .353 & 3/3\\
Regression4096$\times$64 & 9.414 & 3.810 & 3/3\\
\bottomrule
\end{tabular}
\end{center}

Contemporaneous direct controls were $.911$ seconds for Gaussian32 and
$1.028$ for banana8, versus $.610$ and $.629$ for native TFP. The rewrite
was descriptively faster on these controls, while native Stan remained
lower in elapsed time. Three seeds and unequal process/I/O boundaries do
not establish a general efficiency ranking. The centered-funnel result is
failure of those finite-budget configurations, not proof that standard Stan
tuning cannot sample a centered funnel.

\subsection{Native wrapping and genuine GPU execution}

The native bridge benchmark used four sequential native chains, 500 warmup
and 500 retained transitions, with three alternating-order repetitions.
The following are complete-fit milliseconds, excluding compilation. The
native comparator is measured in the same GPU-enabled process as its bridge.

\begin{center}\small
\begin{tabular}{llrrr}
\toprule
Engine & Target & Native CPU & GPU XLA bridge & Ratio\\
\midrule
Stan & Gaussian10 & 36.092 & 36.921 & 1.0230\\
Stan & Banana10 & 55.638 & 56.683 & 1.0188\\
Stan & Centered funnel10 & 190.519 & 191.331 & 1.0043\\
Walnutpie & Gaussian10 & 18.790 & 19.618 & 1.0441\\
Walnutpie & Banana10 & 30.486 & 31.173 & 1.0225\\
Walnutpie & Centered funnel10 & 36.249 & 37.099 & 1.0235\\
\bottomrule
\end{tabular}
\end{center}

The small observed staging overhead accompanies byte-identical tested
outputs, not a speedup from GPU arithmetic. Both short centered-funnel
fits mixed poorly: maximum R-hat was $1.1681$ for Stan and $1.4260$ for
Walnutpie, with minimum bulk ESS $20.7$ and $8.6$. These are not valid
posterior-efficiency comparisons.

The genuine Stan TF/GPU programme retained the shared direct transition and
changed execution structure: static per-chain RNG expansion removed a
serial mapping loop, and command buffers reduced scheduling overhead.
Measured complete-fit seconds for Gaussian10 were $1.4970$ (original GPU
map), $.8253$ (GPU static) and $.5085$ (GPU static with command buffers),
versus $.04475$ for four parallel CmdStan CPU processes. Banana10 gave
$1.7993$, $.9914$, $.6264$ and $.04865$. The cheaper CPU TF Gaussian
versions measured $.07511$ (map) and $.08851$ (static), so a GPU-oriented
layout was not automatically a CPU optimization. GPU Gaussian and banana
posterior screens passed; centered funnel failed in both implementations.
Regression512$\times$32 passed the GPU screen but the short CmdStan run
had R-hat $1.01173$, above the declared $1.01$ threshold. These experiments
verify a generic compiled TF target path and nominate scheduling changes;
they do not establish native-Stan speed or a universal device choice.

\section[Walnutpie validation and execution]{Walnutpie implementation checks and execution improvements}
\label{sec:sw-walnut-validation}

The deterministic port replaced recursion by binary-carry storage without
replacing the source's selection or refinement rules. On each CPU/GPU
backend, 528 native-header comparisons covered 96 transitions, 96 macros,
96 initializations and 240 lifecycle steps on Gaussian, banana, funnel
and regression fixtures. Shared random tapes check discrete work and
random consumption, with scaled FP64 tolerances of $2\times10^{-11}$.
The default lifecycle ran without state resets. Extreme nondefault
target-macro-one tests required complete exposed-input replay after
accumulated roundoff exceeded the trajectory tolerance in two attempts.
Those failures remain recorded: local replay agreement is not unrestricted
long-trajectory identity across backends.

The later initialization/control/R2P qualification passed 48 focused tests,
144 identity-metric author macro fixtures, 144 diagonal conjugacy fixtures
and 216 whole-tree comparisons across three layouts. Deterministic
528-record regressions passed again on CPU and GPU. Whole-tree R2P
comparisons used identity mass; diagonal macro conjugacy alone does not
prove a separate nonidentity whole-tree oracle. Supplied-mass work added
38 CPU regression checks and eight complete GPU cases covering
deterministic/R2P, diagonal/dense and fixed/adaptive combinations.

\subsection{Native C++ versus XLA workloads}

The initial port used four chains and matched starts/data/default settings;
Gaussian10, banana10 and centered funnel10 had 500 warmup plus 500 retained
transitions, while regressions had 100 plus 100. Medians of seven
synchronized fits exclude compilation and reporting.

\begin{center}\small
\begin{tabular}{lrrrr}
\toprule
Target & CPU XLA & GPU XLA & C++ serial & C++ four threads\\
 & \multicolumn{4}{c}{Seconds}\\
\midrule
Gaussian10 & .1206 & .8735 & .0240 & .00849\\
Banana10 & .0871 & .9707 & .0191 & .00869\\
Centered funnel10 & .2028 & 2.1336 & .0322 & .01184\\
Regression4096$\times$128 & 2.3639 & .8234 & 1.7273 & .6574\\
Regression32768$\times$128 & --- & 2.8324 & 29.7394 & 18.3214\\
\bottomrule
\end{tabular}
\end{center}

TF/native batches were not interleaved and their RNGs differ. The large
regression deliberately recomputes its full design matrix at each score.
Its observed GPU advantage over this Eigen/GCC build is therefore a
workload result, not proof of faster Gaussian-regression inference: exact
sufficient-statistic preprocessing changes the comparison substantially.

For the declared quadratic fixture,
\begin{align}
 \log\pi(q)&=-\tfrac12q^\top q-\frac{(Xq-y)^\top(Xq-y)}{2n}\\
 &=-\tfrac12q^\top q-\tfrac12q^\top Gq+q^\top b-\tfrac12c,
 \qquad G=X^\top X/n,\ b=X^\top y/n,\ c=y^\top y/n.
\end{align}
The gradient is $-q-Gq+b$. Both C++ and TF arms implemented this same
algebra, including the constant and nonzero-response term. In a later
separate comparison, full-regression fit medians were $.53280$ seconds on GPU, $.70561$
for generic C++ and $.50423$ for architecture-optimized C++. The
sufficient-statistic medians in the same order were $.40908$, $.02879$
and $.01952$. A nonlinear logistic fixture gave $.50235$, $1.01176$
and $.82944$ seconds. The latter's mean log likelihood plus unit prior
is a tempered synthetic workload, not a full-data posterior claim.

\subsection{Packed storage, failed flattening, and batching}

Packed storage combines twelve tree-span fields into one tensor of shape
$[\text{depth}+1,C,8d+4]$. It reduces stack updates without changing
field values or binary-carry order. The flat variant also replaces nested
leaf/carry traversal by one scheduler. Flat passed numerical checks but
failed its performance criterion on Gaussian and banana; fewer loops
were not sufficient for lower runtime.

The corrected paired experiment used seven seed blocks and balanced AB/BA
orders, four repetitions per seed, four chains and 100 warmup plus 100
retained transitions. Command buffers were enabled equally. Ratios are
geometric means of block-median paired log ratios, with approximate 95\%
Student-$t$ intervals over blocks.

\begin{center}\small
\begin{tabular}{lrrl}
\toprule
Target & Fieldwise s & Packed s & Fieldwise/packed ratio [interval]\\
\midrule
Gaussian10 & .18283 & .17216 & 1.0730 [1.0704, 1.0757]\\
Banana10 & .20934 & .19744 & 1.0681 [1.0656, 1.0706]\\
Centered funnel10 & .36861 & .35926 & 1.0536 [1.0467, 1.0606]\\
Regression4096$\times$128 & .84282 & .81803 & 1.0302 [1.0295, 1.0309]\\
Logistic4096$\times$128 & .74596 & .69624 & 1.0322 [1.0296, 1.0349]\\
\bottomrule
\end{tabular}
\end{center}

An earlier harness retained both paired outputs and used unbalanced order;
its second call was slower regardless of layout. That attempt is excluded.
The corrected experiment releases timed outputs and compares correctness
before timing. Its intervals support a local execution improvement for
packed storage, not an ESS/second or universal performance claim.

Increasing Gaussian chains from 1 to 4 to 16 raised observed throughput
from 1651 to 4886 to 9408 chain-transitions/second, while latency rose
from $.1212$ to $.1637$ to $.3401$ seconds. First-call costs rose from
$6.81$ to $17.56$ to $114.14$ seconds. Static per-chain random generation
increases graph size; more chains are not free parallelism.

\section[Centered funnel replication]{Centered funnel: replication rather than extrapolation}
\label{sec:sw-funnel}

The author funnel with parameter $d=10$ has eleven total coordinates:
$v\sim N(0,9)$ and $x_i\mid v\sim N(0,e^v)$ for $i=1,\ldots,10$.
Earlier local ``funnel10'' fixtures had ten total coordinates. Let
$m$ denote the number of conditional coordinates. Up to the constant
normalization, the shared analytic evaluator is
\begin{align}
 \log\pi(v,x)&=-\frac{v^2}{18}-\frac m2v
                     -\frac12e^{-v}\sum_i x_i^2,\label{eq:sw-funnel-density}\\
 \partial_v\log\pi&=-v/9-m/2+\tfrac12e^{-v}\sum_i x_i^2,
 \qquad \partial_{x_i}\log\pi=-e^{-v}x_i.
\end{align}
Expanding the normal densities gives these formulas and the constant
$-\log3-(m+1)\log(2\pi)/2$. The Python reference includes this
constant; TF omits it, leaving scores and transition probabilities
unchanged. Density/gradient parity was checked against the original SciPy
evaluator. Algebraic equality does not imply bitwise adaptive trajectories:
floating-point changes alter subsequent random decisions and tuning.

The published Figure 13 uses 200,000 iterations and a fixed-step NUTS
comparator selected to approximately match gradient cost. Figure 14 includes
a million-iteration diagnostic. The paper explicitly reports slow mixing
for both methods while WALNUTS-R2P reaches the far neck. It does not assert
that 500 or 2000 retained draws under every Stan configuration must pass
our multi-chain precision screen. Nor are the paper's R2P settings the
deterministic C++ Walnutpie variant used in the early bridge.

\subsection{Author protocol and short replication}

Nine original author files and their MIT licence were preserved unchanged.
The reference adapter keeps the script's normal-with-mean-one start,
1000-iteration Appendix C quantile tuning, last reported tuning parameters
and restart from the original position. Seed zero gives $v=8.292$ in
that start. Analytic tuning of the eleven-coordinate target produced
$h=.3228754012$ and $\delta=.1544232653$, not the paper's rounded
$.36/.21$. Identity mass, maximum depth 12, R2P probability $2/3$,
jitter $.2$ and cap 10 were shared across the comparison arms.

Four chains, 512 archived/excluded settling draws and 8192 retained draws
gave the following short-run results. All rows failed the combined screen,
despite finite states, movement and successful neck visitation.

\begin{center}\small
\begin{tabular}{lrrr}
\toprule
Eleven-coordinate implementation & Max R-hat & Min bulk ESS & Min tail ESS\\
\midrule
Author original & 1.0192 & 308.6 & 501.5\\
Author corrected & 1.0119 & 287.6 & 261.9\\
TF GPU XLA & 1.0080 & 368.8 & 540.2\\
TF CPU XLA & 1.0080 & 368.8 & 540.2\\
\bottomrule
\end{tabular}
\end{center}

The ten-coordinate corrected-author/TF-GPU runs also failed, with minimum
bulk ESS $329.7/319.6$. The GPU paths reached $v=-9.13$ in eleven
coordinates and $-10.87$ in ten. Tail visitation supports the local
refinement mechanism; it does not establish stationary tail probabilities.

\subsection{Completed million-iteration comparison}

The first literal-SciPy campaign was stopped at the owner's request; its
GPU arm separately recorded a CUDA launch failure. Partial evidence was
preserved and not pooled with the restart. A five-repeat one-state
diagnostic measured $153.48\,\mu$s for the SciPy evaluator versus
$7.07\,\mu$s for the analytic one. This explains evaluator overhead,
not a $21.71$-fold whole-chain prediction. Fresh analytic tuning was used
for all four restarted methods, launched concurrently in separate processes.

The final analytic campaign completed one million iterations per chain,
four chains per method. Its primary author-style analysis discards zero
sampling iterations after the tuning restart; exclusion of the first
512 is a separately saved sensitivity. The reported diagnostic screen
includes analytically whitened observables, R-hat and bulk/tail ESS.

\begin{center}\small
\begin{tabular}{lrrrr}
\toprule
Method & Max R-hat & Min bulk ESS & Min tail ESS & Sampling s\\
\midrule
Author original & 1.000254 & 37744.8 & 47093.8 & 4805.1\\
Author corrected & 1.000073 & 36025.6 & 46880.9 & 4894.1\\
TF CPU XLA & 1.000134 & 36259.0 & 45221.5 & 1698.3\\
TF GPU XLA & 1.000145 & 36074.8 & 44373.4 & 15667.9\\
\bottomrule
\end{tabular}
\end{center}

All four passed the within-stream diagnostic screen. Nevertheless, the
saved simultaneous comparison against analytic expectations detects three
eligible differences in the original-author arm: mean $z_0$,
$P(v<-3)$ and $P(v<-6)$. The latter differences are $-.0120075$
with interval $[-.0175579,-.0064571]$ and $-.00351163$ with interval
$[-.00550198,-.00152128]$. The corrected-author arm has one flagged
whitened-coordinate mean, $z_2$, with difference $.00186362$ and
interval $[.00010554,.00362169]$. Neither TF arm has an eligible flagged
contrast in the saved family. The recorded normal threshold is $3.455413$
for family size 91; the intervals use batch MCSE assumptions, with 50/100
block sensitivities preserved.

Thus ``all four are equivalent'' is unsupported even though R-hat and
ESS look satisfactory. The original backward-weight omission is a concrete
implementation concern, but this one comparison does not prove it causes
every discrepancy. Original/corrected author runs share seeds, so their
independent-SE pairwise intervals are explicitly ineligible for inference.
The TF findings support compatibility under the checked observables and
precision, not formal equality of stationary laws. Concurrent execution
also prevents interpreting the runtime column as an isolated speed ranking.

\section{CPU tree control with GPU models}
\label{sec:sw-placement}

The hybrid experiment used a nonlinear regression workload with 16,384
observations, 128 latent bias coordinates, two width-128 tanh stages and
four chains. Fixed generated weights are model data; this is not sampling
all neural-network weights. FP64, TF32 off, identity mass and frozen
$\epsilon=.2$ were common. The hybrid kept a resumable tree controller
on CPU/XLA and complete R2P integration/refinement leaves on GPU/XLA.
Python scheduled the compiled segments and tensor transfers. It was not
one fused heterogeneous executable.

Seven paired seed blocks, two repeats per block and rotated execution order
gave medians of $18.154$ seconds for CPU/XLA with four threads, $1.826$
for full GPU/XLA, and $2.063$ for the hybrid, per two transitions in each
of four chains. The paired hybrid/GPU ratio was $1.13812$ with approximate
95\% interval $[1.12254,1.15392]$. A separate eight-thread CPU sensitivity
gave $10.633$ seconds. Full GPU was the supported execution choice for
this workload; moving control to CPU did not improve speed.

Each hybrid block made thirty leaf round trips carrying about $2.12$ MB
in total, while model tensors stayed resident. Separate raw-copy and
compiled-echo diagnostics took $11.70$ and $133.55$ ms, against an
approximately $237$ ms hybrid/full-GPU difference. Dispatch and
synchronization therefore matter in addition to byte volume; these
measurements are not an additive causal decomposition. All fourteen
primary blocks agreed in discrete work/RNG outputs, with maximum checked
floating difference $1.82\times10^{-12}$.

The later real q20 placement experiment used deterministic Walnutpie and
its frozen NeuTra/filter target, four chains, four seed/start blocks and
balanced execution order. It measured 80 synchronized invocations, each
two transitions per chain, including boundary transfers and output reads.

\begin{center}\small
\begin{tabular}{>{\raggedright\arraybackslash}p{.40\linewidth}r>{\raggedright\arraybackslash}p{.35\linewidth}}
\toprule
Placement & Median s & Paired time/full GPU [95\% interval]\\
\midrule
Full CPU/XLA & 57.529 & 11.3178 [11.0458, 11.5965]\\
Full GPU/XLA & 5.047 & Reference\\
CPU graph tree; GPU target per call & 5.175 & 1.0242 [1.0105, 1.0381]\\
CPU graph tree; complete GPU macro & 5.074 & 1.0057 [.9997, 1.0118]\\
Python tree; complete GPU macro & 5.177 & 1.0273 [.9961, 1.0596]\\
\bottomrule
\end{tabular}
\end{center}

The q20 hybrids use CPU TensorFlow graph/eager control, not the earlier
resumable CPU/XLA tree. These designs must not be conflated. The corrected
benchmark verified all 729 frozen closure hashes, trajectories, work and
RNG counters. Its first attempt was excluded because nominal CPU sampler
calls dispatched to GPU and hybrid RNG lowering differed. Explicit call-site
placement and a shared compiled random primitive repaired that harness.
The final CPU/GPU advantage is supported locally; neither macro hybrid
is distinguishable from full GPU at this precision. These measurements
do not establish the same speed ratio for full80.

\section{Applications through frozen NeuTra transforms}
\label{sec:sw-applications}

\subsection{q20: completed execution without a qualified estimate}

Placement success did not solve the q20 posterior problem. Subsequent
controller repairs distinguished source macro-quality tuning from endpoint
acceptance, preserved every frozen output including repeats, and continued
the same state/RNG without silently retuning. The q20 owner's later
correction made ESS diagnostic-only; its direct floor was removed, while
R-hat, mean MCSE and negative-region-probability precision remained separate
criteria. This q20 policy change does not retrospectively change the
predeclared full80 comparison below.

The final q20 attempt stopped normally at its 10,000-retained-draw limit
per chain on 31 October, with 40,000 total draws. Maximum R-hat was
$1.0816781$ against $1.01$; maximum mean MCSE/SD was $.1542893$
against $.05$; negative-region probability MCSE was $.0469229$ against
$.025$. The prior 8192-draw R-hat was $1.0723291$, so descriptive
progress was not monotone. No q20 posterior estimate is promoted.
This is insufficient mixing/precision in the bounded run, not evidence
that the method or target is mathematically impossible.

\subsection{Full80: compatibility with the saved fixed-HMC posterior}

The requested MacroFinance ``BGL80'' case was resolved from the saved binding
to DZ5 financial full80, not the separate BGL135 model. The frozen
update5000 NeuTra map, data, physical parameter definitions and analytical
score were preserved; no retraining occurred. Four FP64 GPU chains used
TF32 off and exact identity mass in latent coordinates. The saved HMC
reference used $L=5$, $\epsilon=.3363585661$, 2000 excluded warmup
transitions and 1400 retained draws per chain.

An initial $.70$ Walnutpie pilot passed target/cache/mass checks but reached
its 1700-second cap while compiling retained sampling. Its two-transition
warm chunk cost $115.823$ seconds. That pilot was not resumed: source
review corrected the tuning target to $.80$ base-macro quality, initial
step $.1$ plus search, and 50--1000 adaptation transitions checked every
five. The granted additional budget was 70 GPU-worker hours. The corrected
run froze after 105 transitions when relative step agreement reached
$.08708$, below $.10$. Each chain retained its own step size and RNG.
Tuning agreement was not treated as stationarity.

The first posterior assessment at 1400 retained draws per chain passed.
The worker finished at 12:46 Shanghai on 31 October 2026, with supervisor
exit code zero. All warmup draws were excluded; all frozen outputs,
including repeated states, were retained.

\begin{center}\small
\begin{tabular}{lrrl}
\toprule
Diagnostic & Saved HMC & Walnutpie & Required\\
\midrule
Maximum modern R-hat & 1.006456 & 1.005552 & $\leq1.01$\\
Minimum bulk ESS & 948.49 & 1435.83 & $\geq400$\\
Minimum tail ESS & 423.81 & 953.43 & $\geq400$\\
Maximum controller MCSE/SD & .03506 & .02775 & $\leq.05$\\
Recorded worker hours & 5.823 & 20.076 & Descriptive\\
\bottomrule
\end{tabular}
\end{center}

For each physical mean, the comparison formed
\begin{equation}
 Z_j=\frac{\widehat\mu_{W,j}-\widehat\mu_{H,j}}
 {\sqrt{\widehat{\mathrm{MCSE}}_{W,j}^{2}+
             \widehat{\mathrm{MCSE}}_{H,j}^{2}}}.
 \label{eq:sw-mean-comparison}
\end{equation}
Both streams used the same reporting MCSE estimator, distinct from the
controller's precision estimator. The two-sided Bonferroni normal
threshold for 80 means and familywise $.05$ was $3.4205267$. No mean
exceeded it; the largest absolute statistic was $3.1512966$ for
\texttt{US\_Phi\_slope\_level}. The result is compatibility at available
precision, not equality of distributions or proof against shared missed modes.

Warmup and retained sampling each executed GPU/XLA with one qualified trace
and no forbidden host callback. Every chunk passed endpoint cache and
identity-mass checks; retained tuning stayed frozen. GPU allocation growth
was verified. Peak TensorFlow allocator usage was $132.19$ MiB, while
peak host RSS was $15.64$ GiB, reflecting substantial compilation/runtime
overhead outside tensor storage. Excluding the first retained chunk, mean
cost was $46.365$ seconds per transition across four chains. Useful
evaluation counters averaged $16.241$ per chain-transition; inactive
batched target work is not included in that count.

The supervisor charged $72285.101$ seconds, or $20.079$ hours, leaving
$49.921$ hours of the additional allowance unused. The earlier qualification
remains a separate charge. The observed worker-time ratio to HMC is
$3.447$, but HMC's recovery replayed 500 old warmup draws and the methods
have different adaptation/compilation work. This is descriptive accounting,
not a controlled speed ratio. Higher observed Walnutpie ESS does not by
itself establish better ESS per second.

A fresh saved-draw reconstruction reproduced the diagnostics exactly,
verifying ten source snapshots and 2409 sample/checkpoint file hashes.
Its first attempt was blocked by MacroFinance source edits made after
the target graphs had been compiled. The successful review imported only
the checksum-verified pinned diagnostic runtime and never reevaluated the
changed model. The original process retained its compiled functions, but
a fresh target replay from the changed worktree is not reproducible
without recovering the bound source closure. This limitation is recorded
rather than erased by changing expected hashes.

\section{Using and reproducing the implementations}
\label{sec:sw-reproduce}

\subsection{A fixed identity metric for a TensorFlow model}

The following small example illustrates the explicit tensor interface. Its
short fixed counts are a mechanics example, not a posterior protocol. The
Gaussian log density can be replaced by an independent-row TensorFlow model
or by the frozen transformed density in equation~\eqref{eq:sw-neutra-target}.
The factory supplies stable signatures and XLA compilation; the caller must
choose a device, keep data placement explicit, and check validity. GPU use
requires a compatible installed TensorFlow/CUDA environment.

\begin{lstlisting}[language=Python,basicstyle=\ttfamily\scriptsize,breaklines=true]
import os
os.environ['TF_FORCE_GPU_ALLOW_GROWTH'] = 'true'
import tensorflow as tf
gpus = tf.config.list_physical_devices('GPU')
if not gpus:
    raise RuntimeError('This example requires a GPU')
for gpu in gpus:
    tf.config.experimental.set_memory_growth(gpu, True)
    assert tf.config.experimental.get_memory_growth(gpu)
from docs.benchmarks.walnutpie_gpu_tf import (
    Config, make_supplied_mass_functions)

def log_density(q, location):
    return -.5 * tf.reduce_sum((q-location)**2, axis=-1)

cfg = Config(chains=4, dimension=10, warmup=20, draws=20,
             adapt_mass=False)
f = make_supplied_mass_functions(tf, log_density, cfg,
    model_signature=[tf.TensorSpec([10], tf.float64)])
with tf.device('/GPU:0'):
    q = tf.ones([4, 10], tf.float64) * .2
    keys = tf.constant([[101, 0], [102, 0],
                        [103, 0], [104, 0]], tf.int32)
    mass = tf.ones_like(q)
    epsilon = tf.fill([4], tf.constant(.1, tf.float64))
    fit = f.fit(q, keys, mass, epsilon,
                tf.zeros([10], tf.float64))
assert bool(tf.reduce_all(fit.valid).numpy())
tf.debugging.assert_equal(fit.warmup.state.inverse_mass, mass)
retained = fit.sampling.draws  # [20, 4, 10]; warmup excluded
\end{lstlisting}

Here \texttt{adapt\_mass=False} freezes the supplied momentum covariance,
but step-size and minimum-work adaptation still operate during warmup.
The factory's \texttt{adapt\_step=False} option skips only the initial
step search; it does not disable warmup adaptation. To use already frozen
R2P tuning, set zero warmup, provide the step explicitly, skip that search,
and invoke sampling on the initialized state. A change of model data or
target requires reinitialization of density/score caches and reassessment
of tuning. Continuing an unchanged sampling state preserves the RNG and
frozen tuning.

The numerical state contains position, density, score, random keys and
counters, step size, inverse mass, minimum work, adaptation iteration,
Adam accumulators and discounted moments. Dense initialization additionally
stores the fixed basis. Run-control checkpoints also contain iteration
counts and log-density moments. Save this complete state for continuation;
saving only the last position does not reproduce the run. Fixed-capacity
chunk outputs carry a mask, and padded rows must never enter estimates.

\subsection{Tuning and comparison protocol}

For the standalone Stan-style route, first validate the unconstrained
target and score, Jacobians, row isolation, finite starts and metric
convention. Use the inspected Stan three-stage warmup, its metric-boundary
step search and dual-averaging restart, then freeze adaptation before
retaining samples. The adaptation-only wrapper rejects unsupported short
budgets; the direct port separately preserves the pinned Stan short-budget
fallback described in Section~\ref{sec:sw-stan-tuning}. These are different
interface contracts. Source defaults nominate a starting configuration;
they do not certify that a target is adequately explored.

For Walnutpie, record the deterministic or R2P variant, base interval,
refinement bounds, jitter, acceptance statistic, mass convention and
initialization search limit. Its default Adam warmup is not Stan dual
averaging, and R2P combined with that warmup is not the paper's Appendix C
quantile algorithm. The author-replication driver runs Appendix C in the
independent Python reference and transfers its frozen parameters into XLA.
An upstream log-density stopping statistic or step/mass agreement can
finish adaptation without establishing parameter convergence.

Compare tuned fixed-length HMC with a declared search over step size and
leapfrog count, charging its calibration and selection work. Compare
NUTS with its complete adaptation and charge that work too. Use held-out
retained streams for the posterior comparison. Report cold compilation,
warmup, retained execution, diagnostics and total elapsed time separately;
give the exact number of chains, device, dtype, model/gradient backend and
work-counter meaning. For execution-layout comparisons, use common inputs,
balanced order, synchronization and uncertainty over independent seed blocks.
For different samplers, assess validity and posterior agreement before
ranking effective information per unit time. Shared seeds do not justify
an independent-stream standard-error formula.

\subsection{Source and test map}

The authoritative implementation and test locations, relative to the
repository root, are listed below. Reported historical test totals refer
to their recorded revisions, not a promise that a future suite has the
same test count. Native-header tests require the pinned local sources;
missing dependencies must not silently substitute a different reference.

{\small
\begin{longtable}{>{\raggedright\arraybackslash}p{.24\linewidth}>{\raggedright\arraybackslash}p{.68\linewidth}}
\toprule
Component & Implementation and evidence\\
\midrule
\endhead
Stan warmup & \path{docs/benchmarks/stan_nuts_warmup.py};
\path{docs/reference/nuts-tuning.md}; adaptation code-audit result.\\
Direct Stan and TFP & \path{docs/benchmarks/stan_nuts.py},
\path{docs/benchmarks/stan_nuts_tfp.py};
\path{tests/test_stan_nuts_tfp.py}; direct-port and native-TFP result notes.\\
Native bridge & \path{docs/benchmarks/native_sampler_bridge/};
native Stan/Walnutpie result note and source manifests.\\
Stan GPU models & \path{docs/benchmarks/stan_gpu_tf/};
\path{tests/test_stan_gpu_tf.py}.\\
Walnutpie numerical core & \path{docs/benchmarks/walnutpie_gpu_tf/core.py},
\path{api.py}, \path{r2p.py} in that package;
\path{tests/test_walnutpie_gpu_tf.py}.\\
Initialization, control, mass & \path{initialization.py},
\path{control.py}, \path{mass_coordinates.py} in the same package;
\path{tests/test_walnutpie_gap_closure.py},
\path{tests/test_walnutpie_supplied_mass.py}.\\
Funnel and hybrid studies & \path{author_reference/},
\path{funnel_master.py}, \path{replication.py},
\path{replication_analysis.py}, \path{hybrid_benchmark.py} in that package;
\path{tests/test_walnutpie_hybrid.py}.\\
Full80 application & \path{docs/benchmarks/full80_walnutpie_sampling.py},
its support, controller-helper and report modules;
\path{tests/test_full80_walnutpie.py}.\\
\bottomrule
\end{longtable}}

The Walnutpie C++ oracle is also saved under
\path{docs/benchmarks/walnutpie_gpu_tf/prebuilt/linux-x86_64/20261029/}.
Its executable is 249,008 bytes (243.2 KiB); the matching source archive
is about 1 MiB. Checksums, compiler flags, pinned headers, Eigen 3.4.0 and
licences accompany it. It requires Linux x86-64, glibc 2.34 or newer and
libstdc++ symbols through \texttt{GLIBCXX\_3.4.29}/\texttt{CXXABI\_1.3.9}.
It has no TensorFlow or CUDA dependency. It emits 528 deterministic JSON
oracle records and implements the supplied analytic targets, not arbitrary
user TensorFlow models. The recorded generic build took 7.33 seconds.
Platform-specific compilation remains necessary elsewhere, and each new
TensorFlow model signature still requires XLA compilation.

For a bounded CPU correctness check with the prebuilt reference, use a fresh
output directory and the installed TensorFlow environment:
\begin{lstlisting}[basicstyle=\ttfamily\scriptsize,breaklines=true]
CUDA_VISIBLE_DEVICES=-1 TF_FORCE_GPU_ALLOW_GROWTH=true \
python -m docs.benchmarks.walnutpie_gpu_tf.worker \
  --mode validate --device cpu --execution-layout packed \
  --oracle docs/benchmarks/walnutpie_gpu_tf/prebuilt/linux-x86_64/20261029/oracle \
  --output /tmp/walnutpie-validation-new

CUDA_VISIBLE_DEVICES=-1 TF_FORCE_GPU_ALLOW_GROWTH=true \
python -m pytest -q tests/test_full80_walnutpie.py
\end{lstlisting}

Do not pass the tracked executable as a build destination: the original
master's build, qualify and all phases rebuild that path. Worker validation
or the master's validate/benchmark phase can consume the prebuilt directly.
Full GPU benchmarks and million-iteration replications are separate,
budgeted experiments; their launch commands, source pins and stop conditions
belong to the saved plans, rather than this small correctness example.

\subsection{Result records and reproducibility limits}

The result-note index below gives filenames under \path{docs/plans/}.
Each note links its run manifests, exact commands, seeds, environment,
timing boundaries and structured diagnostics. The monograph evidence
snapshot under \path{docs/reference/stan-walnutpie-evidence-20261031/}
preserves selected completed-result JSON files, their original locations
and SHA-256 hashes. Large raw draws, machine-local source installations
and financial-model inputs remain at the recorded experiment locations;
the compact snapshot is not a replacement for those inputs.

{\small
\begin{longtable}{>{\raggedright\arraybackslash}p{.23\linewidth}>{\raggedright\arraybackslash}p{.69\linewidth}}
\toprule
Study & Result note or completed record\\
\midrule
\endhead
Early devices/HMC & \path{bayesfilter-nuts-xla-benchmark-result-2026-10-27.md};
\path{bayesfilter-nuts-fixed-hmc-result-2026-10-27.md};
\path{bayesfilter-nuts-geometry-overhead-result-2026-10-27.md}.\\
Stan transfer/audit & \path{bayesfilter-stan-nuts-transplant-result-2026-10-27.md};
\path{bayesfilter-stan-adaptation-code-audit-result-2026-10-27.md};
\path{bayesfilter-full-stan-nuts-port-result-2026-10-27.md}.\\
TFP and GPU Stan & \path{bayesfilter-stan-tfp-native-result-2026-10-28.md};
\path{bayesfilter-native-stan-walnutpie-result-2026-10-28.md};
\path{bayesfilter-stan-gpu-tf-result-2026-10-28.md}.\\
Walnutpie transfer & \path{bayesfilter-walnutpie-gpu-tf-result-2026-10-28.md};
\path{bayesfilter-walnutpie-gap-closure-result-2026-10-29.md};
\path{bayesfilter-walnutpie-supplied-mass-result-2026-10-29.md}.\\
Optimization/hybrid & \path{bayesfilter-walnutpie-optimization-result-2026-10-29.md};
\path{bayesfilter-walnutpie-hybrid-large-model-result-2026-10-29.md}.\\
Author funnel & \path{bayesfilter-walnuts-author-funnel-result-2026-10-29.md};
million-run \path{REPORT.md} and \path{comparison.json} in
\path{docs/benchmarks/artifacts/walnuts-million-analytic-20261029-01/}.\\
q20 & \path{bayesfilter-q20-walnutpie-placement-results-2026-10-29.md};
\path{bayesfilter-q20-walnutpie-estimation-results-2026-10-29.md},
including its 31 October terminal update.\\
Full80 & \path{bayesfilter-full80-walnutpie-result-2026-10-30.md};
\path{comparison-01/} and \path{terminal-review-02/} under
\path{docs/benchmarks/artifacts/full80-walnutpie-20261030/}.\\
\bottomrule
\end{longtable}}

The remaining limits are specific. Appendix C tuning has not been
transferred to a native TensorFlow kernel. The exact extreme-start
$v=-30$ transient study is incomplete. Dense supplied-mass adaptation is
diagonal in a fixed basis. The C++ bridge does not fuse arbitrary TF
models, and the pure TensorFlow port does not reproduce upstream random
streams bit for bit. CPU/GPU preference depends on model cost, batching
and compiler behaviour. The full80 saved-draw comparison is complete, but
fresh target replay requires recovery of its bound MacroFinance source
closure; q20 still lacks an accepted posterior estimate. These distinctions
identify what further evidence is needed without discarding the completed
source checks, execution measurements or posterior comparisons.


\clearpage
\endgroup
